\section{Numerical Experiments and Forecasting Boundaries}

\subsection{The retained production benchmark}

The current D-Link JSON benchmark compares both backends on the same
846-by-4 fit matrix, 288 calibration windows, and 270 development windows.
It builds temporary artefacts rather than replacing the accepted model.
Complete fit timing excludes data loading and saving. Likelihood timing uses
200 calls per batch and score timing 1,000, with five batches and warm-up.
The fit median uses five measurements.

\begin{table}[H]
\centering\small
\begin{tabular}{lrrrr}
\toprule
\textbf{Operation} & \textbf{LM} & \textbf{SciPy} & \textbf{LM/SciPy} & \textbf{Unit} \\
\midrule
Complete fit & 192.974 & 181.286 & 1.064 & ms \\
Fit-matrix kernel & 55.457 & 49.961 & 1.110 & $\mu$s \\
Single-window score & 12.044 & 13.394 & 0.899 & $\mu$s \\
\bottomrule
\end{tabular}
\caption{Exact timings from the retained D-Link production benchmark JSON, rounded for display.}
\label{tab:retained-backend}
\end{table}


The retained environment is Python 3.12.14, NumPy 2.5.1, and SciPy 1.18.0 on
Linux x86-64. These are the versions in that record, not the versions pinned
for the later edge image. The separately retained three-dataset figure belongs
to the earlier report's timing run. Neither is an observation that the entire
software was freshly benchmarked during this documentation update.

The numerical comparison finds a maximum relative alpha difference of
$3.8571\times10^{-7}$ and maximum absolute score difference of
$1.0917\times10^{-5}$, with zero decision mismatches in 270 windows.
Both backends flag three windows. These small differences establish agreement
on this workload; they do not prove every possible threshold decision is
insensitive to numerical error.

\subsection{Why the matrix optimisation matters}

The authors' numerical source was preserved, while a separate production copy
was adapted for safe Python use. The older adapter repeated allocation and the
adaptive Euler--Maclaurin/horizontal-shift setup for every row.
That setup depends on alpha, category count, and requested precision, not on
each row's observed counts. Moving it outside the row loop removes repeated
work while preserving the sum of independent-row kernels.

The recorded native experiment compared old and optimised libraries in 1,000
seeded cases and found bit-identical results. It reduced the D-Link LM matrix
call from 114.68 to 55.46 microseconds and complete fit time from 0.30410 to
0.19297 seconds. The retained SciPy fit is 0.18129 seconds. The useful claim is
that an identified integration cost was removed; it is not that LM always
outperforms SciPy in this production workload.

\subsection{Paper-shaped aggregate experiment}

\path{scripts/lm_professor_experiment.py} separately follows the supplied
study's aggregate convergence architecture. It uses seven D-Link days,
10/15/20/25-percent samples, three fixed repeats, moment-based initialisation,
row-wise digamma updates, and a pooled-vector stopping likelihood.
The two backends receive identical samples.

There are 84 paired fits, of which 79 converge for both backends.
The five matching failures occur in small samples approaching the multinomial
boundary. Across all pairs, recorded likelihood time falls from 24.3311 to
21.3625 seconds and wall time from 33.7594 to 30.7930 seconds with LM.
These correspond to reductions of 12.2\% and 8.8\%.

Local daily row counts differ from those in the source paper, and its precise
sampling realisations are not recoverable. The experiment is therefore a
controlled architectural/backend comparison on local data, not a literal
reproduction of every published table. A pooled-vector likelihood is also not
equal to a sum of row likelihoods. Its speed cannot justify silently replacing
the production objective.

\subsection{Where numerical accuracy becomes important}

The retained high-precision probe compares aggregate kernels with a
100-decimal-digit mpmath reference. At the current concentration, about 122,
both backends are already accurate. At concentration $10^{10}$, the recorded
absolute errors are approximately $2.83\times10^{-5}$ for SciPy and
$5.65\times10^{-8}$ for LM; at $10^{14}$ they are about $0.337$ and $0.00103$.
Subtracting large, close log-Gamma values is increasingly difficult in that
regime. These are specific probes of numerical behaviour, not measured
improvements in attack recall.

The native wrapper checks library availability and coefficient loading, converts
inputs to contiguous doubles, and serialises calls around the supplied global
state. Non-finite precision failures return typed Python errors. Sparse, zero,
large-concentration, repeated-call, and backend agreement tests exercise this
boundary. The original mathematical approximation and licence are retained.

\subsection{Two different forecasting paths}

The installed \texttt{forecast} command is a scaffold. It draws synthetic
multinomial outcomes from the same fitted mean that supplies its expected-count
forecast, using a fixed local seed and a total of 100 per hour. Its Brier and
Jensen--Shannon outputs exercise helper code; they are not accuracy measurements
against independently observed future traffic.

The paper-shaped script does perform a narrower real-count comparison:
1,000 Dirichlet samples from each 25-percent daily model are averaged and
compared with observed proportions on October 19, 21, and 22.
Identical random streams make the backend comparison fair. The largest
recorded Brier difference is $2.84\times10^{-13}$, far smaller than the
roughly $10^{-4}$ changes caused by choosing another fixed training sample.
This experiment does not make the installed forecast scaffold an empirical
forecasting service.

\subsection{Small metric utilities}

The evaluation package also contains total variation distance, KL divergence,
mean squared error, Brier score, and Jensen--Shannon divergence.
Their inputs and reduction rules matter: the Brier helper averages the
per-row squared probability error, while the current Jensen--Shannon helper
uses the summed KL terms supplied to it. These utilities support modelling
experiments; their existence is not evidence that the detector has been
benchmarked against Poisson, neural, or countermeasure-selection alternatives.
